\documentclass[11pt]{article}
\usepackage[a4paper,margin=18mm]{geometry}
\usepackage[no-math]{fontspec}
\setmainfont{Latin Modern Roman}
\usepackage{amsmath,amssymb,amsthm,xfrac,xspace,hyperref,graphicx,comment}
\excludecomment{DefTralics}
\providecommand{\backmatter}{}

\providecommand{\bibliofont}{\small}
\newcommand{\ie}{i.e.,\xspace}
\newcommand{\eg}{e.g.,\xspace}
\newcommand{\etc}{etc.\xspace}
\newcommand{\cf}{cf.\xspace}
\newcommand{\resp}{resp.\xspace}
\newcommand{\smaller}{\small}
\newcommand{\Subsection}{\subsection}
\newcommand{\Subsubsection}{\subsubsection}
\newcommand{\skpt}{\smallskip}
\emergencystretch=3em
\input{author-preamble.tex}
\begin{document}
\begin{center}{\Large Stable item 2052}\end{center}
\paragraph{Source and attribution.}
Marcello Bernardara, Enrico Fatighenti, Laurent Manivel, \emph{Nested varieties of K3 type}, Journal de l’École polytechnique — Mathématiques 8 (2021), publisher version, DOI 10.5802/jep.156. \href{https://jep.centre-mersenne.org/articles/10.5802/jep.156/}{Publisher article}. Authors' text: \href{https://creativecommons.org/licenses/by/4.0/}{CC BY 4.0}.
\paragraph{Editorial problem boundary.}
Prove only Theorem 3 below for a smooth complex hyperplane section of \(\Gr(k,V_n)\), with \(n>3k\) and \(k>2\), in the exact author definition of \(N\)-Calabi--Yau type for \(N=k(n-k)+1-2n\). The middle-Hodge vanishing and extremal one-dimensional summand are the selected conclusion. All partition calculations, twisted-form and conormal Euler-characteristic context are retained. The later strong-Calabi--Yau deformation-map suggestion, whose proof the authors explicitly did not check, is outside the selected claim.
\paragraph{Difficulty (editorial): H2.}
The allowed Borel--Bott--Weil theorem supplies a rule for individual twisted-form summands but does not establish the required vanishing range uniformly over all contributing partitions. The author derives the restrictive partition inequalities and classifies the boundary cases, then combines their exact Euler characteristics with the hyperplane conormal formula. Those new calculations identify the first nonzero middle Hodge summand and its dimension without a coprimality assumption.
\paragraph{Editorial transcription note.}
Original author language, mathematics and ordering are preserved. Publisher front matter is replaced by this independent article wrapper. Bibliography is recovered from matching cached publisher HTML, including its TeX formulas. No new proof or mathematical repair is supplied. Surrounding original results are retained for conservative context and are not extra selected corpus claims.
\clearpage
\input{author-body.tex}
\begin{thebibliography}{99}
\bibitem{at12} N. Addington \& R. Thomas - “Hodge theory and derived categories of cubic fourfolds” , Duke Math. J. 163 (2014) no. 10, p. 1885-1927
\bibitem{aw} M. Andreatta \& J. A. Wiśniewski - “A note on nonvanishing and applications” , Duke Math. J. 72 (1993) no. 3, p. 739-755
\bibitem{BLMS} A. Bayer, M. Lahoz, E. Macrì \& P. Stellari - “Stability conditions on Kuznetsov components” , 2017, Appendix joint with X. Zhao
\bibitem{beauville-coble} A. Beauville - “The Coble hypersurfaces” , Comptes Rendus Mathématique 337 (2003) no. 3, p. 189-194
\bibitem{beauville-donagi} A. Beauville \& R. Donagi - “La variété des droites d’une hypersurface cubique de dimension $4$” , C. R. Acad. Sci. Paris Sér. I Math. 301 (1985) no. 14, p. 703-706
\bibitem{vlad-bi} V. Benedetti - “Bisymplectic Grassmannians of planes” , 2018
\bibitem{bondal-orlov-archive} A. I. Bondal \& D. O. Orlov - “Semiorthogonal decomposition for algebraic varieties” , 1995
\bibitem{craw} A. Craw - “An introduction to motivic integration” , in Strings and geometry , Clay Math. Proc. , vol. 3 , American Mathematical Society, Providence, RI, 2004, p. 203-225
\bibitem{dpfmr} P. De Poi, D. Faenzi, E. Mezzetti \& K. Ranestad - “Fano congruences of index 3 and alternating 3-forms” , Ann. Inst. Fourier (Grenoble) 67 (2017) no. 5, p. 2099-2165
\bibitem{debarre-voisin} O. Debarre \& C. Voisin - “Hyper-Kähler fourfolds and Grassmann geometry” , J. reine angew. Math. 649 (2010), p. 63-87
\bibitem{Ein-ShepB} L. Ein \& N. Shepherd-Barron - “Some special Cremona transformations” , Amer. J. Math. 111 (1989) no. 5, p. 783-800
\bibitem{eg2} E. Fatighenti \& G. Mongardi - “A note on a Griffiths-type ring for complete intersections in Grassmannians” , 2018
\bibitem{fm} E. Fatighenti \& G. Mongardi - “Fano varieties of K3 type and IHS manifolds” , Internat. Math. Res. Notices (2021) no. 4, p. 3097-3142
\bibitem{fik-griffiths} D. Favero, A. Iliev \& L. Katzarkov - “On the Griffiths groups of Fano manifolds of Calabi-Yau Hodge type” , Pure Appl. Math. Q 10 (2014) no. 1, p. 1-55
\bibitem{gsw} L. Gruson, S. V. Sam \& J. Weyman - “Moduli of abelian varieties, Vinberg $\theta $-groups, and free resolutions” , in Commutative algebra , Springer, New York, 2013, p. 419-469
\bibitem{hassett:special-cubics} B. Hassett - “Special cubic fourfolds” , Compositio Math. 120 (2000) no. 1, p. 1-23
\bibitem{Huybook} D. Huybrechts - Fourier-Mukai transforms in algebraic geometry , Oxford Math. Monographs , The Clarendon Press, Oxford University Press, Oxford, 2006
\bibitem{huybrechts-rennemo} D. Huybrechts \& J. V. Rennemo - “Hochschild cohomology versus the Jacobian ring and the Torelli theorem for cubic fourfolds” , Algebraic Geom. 6 (2019) no. 1, p. 76-99
\bibitem{imcy} A. Iliev \& L. Manivel - “Fano manifolds of Calabi-Yau Hodge type” , J. Pure Appl. Algebra 219 (2015) no. 6, p. 2225-2244
\bibitem{Jiang-Leung} Q. Jiang \& N. C. Leung - “Derived category of projectivization and flops” , 2018
\bibitem{kuz-lefschetz} A. Kuznetsov - “Lefschetz decompositions and categorical resolutions of singularities” , Selecta Math. (N.S.) 13 (2008) no. 4, p. 661-696
\bibitem{kuz-hochsch} A. Kuznetsov - “Hochschild homology and semiorthogonal decompositions” , 2009
\bibitem{kuz:4fold} A. Kuznetsov - “Derived categories of cubic fourfolds” , in Cohomological and geometric approaches to rationality problems , Progress in Math. , vol. 282 , Birkhäuser Boston, Boston, MA, 2010, p. 219-243
\bibitem{kuz-fractional} A. Kuznetsov - “Calabi-Yau and fractional Calabi-Yau categories” , J. reine angew. Math. 753 (2019), p. 239-267
\bibitem{leuschke} G. J. Leuschke - “Non-commutative crepant resolutions: scenes from categorical geometry” , in Progress in commutative algebra 1 , de Gruyter, Berlin, 2012, p. 293-361
\bibitem{Li-Per-Xiao} C. Li, L. Pertusi \& X. Zhao - “Twisted cubics on cubic fourfolds and stability conditions” , 2018
\bibitem{lunts} V. A. Lunts - “Categorical resolution of singularities” , J. Algebra 323 (2010) no. 10, p. 2977-3003
\bibitem{schubert} L. Manivel - Symmetric functions, Schubert polynomials and degeneracy loci , SMF/AMS Texts and Monographs , vol. 6 , American Mathematical Society, Providence, RI; Société Mathématique de France, Paris, 2001
\bibitem{orlovprojbund} D. O. Orlov - “Projective bundles, monoidal transformations, and derived categories of coherent sheaves” , Izv. Ross. Akad. Nauk Ser. Mat. 56 (1992) no. 4, p. 852-862
\bibitem{rennemo-segal} J. V. Rennemo \& E. Segal - “Hori-mological projective duality” , Duke Math. J. 168 (2019) no. 11, p. 2127-2205
\bibitem{snow} D. M. Snow - “Cohomology of twisted holomorphic forms on Grassmann manifolds and quadric hypersurfaces” , Math. Ann. 276 (1986) no. 1, p. 159-176
\bibitem{sommese} A. J. Sommese - “Submanifolds of Abelian varieties” , Math. Ann. 233 (1978) no. 3, p. 229-256
\bibitem{voisin-book} C. Voisin - Théorie de Hodge et géométrie algébrique complexe , Cours Spécialisés , vol. 10 , Société Mathématique de France, Paris, 2002
\bibitem{weyman-book} J. Weyman - Cohomology of vector bundles and syzygies , Cambridge Tracts in Math. , vol. 149 , Cambridge University Press, Cambridge, 2003
\end{thebibliography}
\end{document}
